% Pursuing Stacks (cotes 134-2 to 134-8): archive cote 134-3, site page 16 (PDF page 17), typescript page 16.
% typescript. Draft compared: ps2.tex:716-774.
% AI-transcribed by Claude Opus 5.5 (claude-opus-5-5, adaptive thinking, effort high), Message Batches API, batch ps1 (msgbatch_01Rv5hCaNTtUt4JHx1yqoHPL), 2026-11-12; not yet checked by eye.
\scanpage{134-3/16}
\mspage{16}
the corresponding "diagonal" complex,  which is just a usual complex and
adequately describes the "homotopy type" of the multicomplex. So why
bother with relatively messy kind of models,  when just usual complexes
suffice !  Here however the point is not to get the handiest possible
model categories (whatever our cr\over{e}{i}teria of "handyness"),  but rather to
get an idea of the variety of algebraic structures suitable for defining
homotopy types,  and perhaps to come to a clue of what is common to all
these.  Moreover,  I feel the relation between ordinary comple\over{c}{x}es,
and n-multicomplexes,  is of much the same nature as the relation bet-%
ween just ordinary categories,  perfectly sufficient for describing homo
-topy types,  and  n-categories  or n-stacks.  This reminds of course of
Quillen's idea of defining n-categories in terms of n-semisimplicial
complexes,  in much the same way as categories can be defined (via the
Nerve functor)  in terms of usual ss complexes.  I hit again upon multi-%
simplices (without symmetry),  when trying to reduce to a \uline{minimum} the
\over{category}{category} $\underline{B}_o$ of \xout{\unsure{socalled}} what should be called "standar\over{t}{d}"amalgamated
sums of the cells $D_n$ ,  where my tendency initially (L, \deleted{\unsure{p.7}}\inserted{section 8}) had rather
been to be as generous as possible,  in order to be as stringent as ima-%
ginable  for the completion condition (A)  of (L, \deleted{\unsure{p.8}}\inserted{section 9}).  Now it turns
out that the coherence relations which seem to have been written down so
far (and \inserted[typed]{the like of} which presumably will suffice to imply full completeness of
coherence relations,  in the sense of (A))  make use only of very restri-%
ted types of such amalgamated sums,  expressible precisely in terms of
multisimplices.  (This I ckecked \inserted[typed]{for instance} on the full list of data and axioms
for Bénabou's "bicategories" namely  2-stacks,  in \hand{his} 1967 Midwest Cate-%
gory Seminar \inserted{exposé} (already referred to.)  I'll have to come back upon this
point with some care,  which gives also a pretty natural way for getting
Quillen's functor from  n-stacks to n-ss complexes.
\margin{[hand bracket in the left margin along the last four lines]}

\margin{22.} \inserted{All} \over{Th}{th}ese examples of possible models for homotopy types can be viewed
as generalizations of usual complexes,  or of usual categories\over{.}{;}  I would
like to give a few others which go in the opposite direction\deleted{s}  -  they
may be viewed as  \uline{particular cases} of categories.  One is the (pre)order
structure,  which may be viewed as a category structure when the
map    Fl $\longrightarrow$ Ob x Ob  defined by the source and target maps is injec-%
tive.  Such category is equivalent \deleted{to} (and hence homotopic) to the
category associated with the corresponding \uline{ordered} set (when $x\geqslant y$ and
$y\geqslant x$ imply $x=y$)\over{,}{.} \deleted{which is} \inserted{Ordered sets are} more familiar I guess as \deleted{a} model object\hand{s} for
describing combinatorially a topological space,  in terms of a "cellular
subdivision"  by \over{\illegible{1}}{compact}\inserted{(17)} subsets or "cells" ("strata" would be a more
\margin{+}
